We give the construction, including the boundary cases of every division.  Choose first the support square \(S=[-1/2,1/2]^2\), the interface \(\{(u,0):-1/2\leq u\leq1/2\}\), and the two half-square assignment regions.  The base density is \(f_*=1\).  This geometry satisfies all chart, reach, length, compactness, and two-sided-thickness requirements uniformly: at an endpoint the smallest relevant intersection is a quarter disk, and for \(r\leq r_0\leq1/4\) its probability is bounded below by a fixed multiple of \(r^2\), larger than the allowed \(c_-r^2\).

Let \(\phi\) be a nonnegative radial \(C^\infty\) bump, supported on the unit disk, equal to one at the origin, and with all derivatives through order \(q\) bounded.  Put
\[
 \Delta=d a_n,\qquad w=(d_1\Delta)^{1/q},
\]
where \(d,d_1>0\) will be chosen sufficiently small using only the displayed class constants.  Select points \(x_1,\ldots,x_M\) on the \(\delta\)-trimmed interface whose mutual distances are at least \(3w\).  Elementary interval packing gives \(M\geq c w^{-1}\geq c_0a_n^{-1/q}\).  Set \(\phi_j(z)=\phi((z-x_j)/w)\).  The supports are disjoint, and for every \(k\leq q\),
\[
 \lVert D^k(\Delta\phi_j)\rVert_\infty
 \leq C_\phi\Delta w^{-k}
 \leq C_\phi d_1^{-1}.
\]
Thus disjoint sums of these bumps remain in the same \(q\)-H\"older ball after increasing the fixed constant \(d_1\) and decreasing the fixed baseline slope below.

Take a small \(b>0\) and define the control mean \(p_0(z)=1/2\) and, for \(\omega\in\{0,1\}^M\), the treated mean
\[
 p_{1,\omega}(z)=\frac12+bz_1+\Delta\sum_{j=1}^M\omega_j\phi_j(z).
\]
The constants may be chosen so that all these means lie in \([1/4,3/4]\) and their H\"older norms are at most \(L\).  Conditional on \(X=z\), let
\[
 Y(t)=B_t+\varepsilon_t,qquad B_t\sim\operatorname{Bernoulli}(p_{t,\omega}(z)),qquad
 \varepsilon_t\sim N(0,1),
\]
with the displayed variables conditionally independent.  Then \(\mu_{t,P_\omega}=p_{t,\omega}\), the conditional variance belongs to \([1,5/4]\subseteq[L^{-1},L]\), and Hoeffding's lemma followed by the Gaussian moment-generating identity gives the required conditional sub-Gaussian inequality with envelope \(L\).

It remains to define the angular density correction.  In polar coordinates \(z=x_j+r(\cos\theta,\sin\theta)\), choose a continuous cutoff \(\chi\) which is zero on \([0,1]\) and one on \([2,\infty)\), and put, on the disk about \(x_j\),
\[
 A_n(r)=-\frac{2\Delta\phi(r/w)}{br}\,\chi\!\left(\frac{br}{8\Delta}\right),
 \qquad
 f_\omega(z)=f_*\left\{1+\omega_jA_n(r)\cos\theta\right\}.
\]
Set the correction to zero outside the disjoint disks.  There is no division at zero because the cutoff vanishes on a neighborhood of zero.  On its active region \(br\geq8\Delta\), \(|A_n(r)|\leq1/4\); it also vanishes continuously at the inner and outer boundaries.  Hence \(f_\omega\) is continuous, integrates to one because \(\int_0^{2\pi}\cos\theta\,d\theta=0\), and lies between \(3/4\) and \(5/4\), which is inside \([L^{-1},L]\).

For either half-circle, \(\int\cos\theta\,d\theta=0\).  Consequently the radial mass integral is exactly the base one for both treatment arms.  On the treated upper half-circle and on radii where the cutoff equals one, radiality of \(\phi\) and
\(\int_0^\pi\cos^2\theta\,d\theta=\pi/2\) give
\[
\begin{aligned}
 &\int_0^\pi p_{1,\omega}(x_j+r e_\theta)
       \{1+A_n(r)\cos\theta\}\,d\theta
 -\int_0^\pi \left(\frac12+b(x_j)_1+br\cos\theta\right)d\theta\\
 &=\pi\Delta\phi(r/w)+brA_n(r)\frac\pi2=0.
\end{aligned}
\]
For the control arm the corresponding equality follows from the constant mean.  A mixture of the variables \(B_t+\varepsilon_t\) over angles depends on the angular law only through the averaged Bernoulli success probability.  Therefore the entire conditional law of \(Y\), not merely its mean, agrees between adjacent vertices at center \(x_j\) outside \(r\leq16\Delta/b\).

Inside that radius the radial densities still agree, all mixture weights stay away from zero and one, and the two Bernoulli success probabilities differ by at most \(C\Delta\).  The elementary inequality
\[
 \operatorname{kl}(u,v)\leq\frac{(u-v)^2}{v(1-v)},\qquad 0<u,v<1,
\]
is obtained by applying \(\log x\leq x-1\) to the two terms of binary relative entropy.  Adding the common Gaussian noise can only decrease relative entropy, by conditioning and Jensen's inequality.  Since the radial density is bounded by \(Cr\), one observation consequently satisfies
\[
 \operatorname{KL}(P_{\omega,x_j}^{\mathrm{dist}},P_{\omega^{(j)},x_j}^{\mathrm{dist}})
 \leq C\int_0^{16\Delta/b}\Delta^2r\,dr\leq C'\Delta^4.
\]
Tensorization gives \(C'n\Delta^4=C'd^4\log n\).  Since
\(\log M=(1/q+o(1))\log(a_n^{-1})=(1/(4q)+o(1))\log n\), choosing \(d\) sufficiently small gives the stated bound with any fixed \(\alpha<1/8\).  At \(x_j\), adjacent target curves differ by exactly \(\Delta\), proving separation with \(c_1=d\).

All class conditions have now been checked.  For \(L>4\), repeat the construction on \([-1,1]^2\) with base density \(1/4\).  The strict slack \(1/4-L^{-1}>0\) permits the relative amplitude of \(A_n\) to be reduced, while retaining the cancellation identity by reducing \(b\) and \(d\) proportionally.  The prescribed chart and assignment regions are then exactly those of \(\mathcal P_{\mathrm{straight}}\).  At \(L=4\) such a square perturbation is impossible because \(f\geq1/4\) on a set of area four and \(\int f=1\) force \(f\equiv1/4\); this proves both the asserted \(L>4\) qualification and the necessity of narrowing the frozen statement.